\chapter{Parametric variables}\label{ch:parameters}

Assignment data is often algebraic: a cable \emph{217 + A} metres long, a load of
\emph{(117 -- A)}~kW. \app\ lets you write exactly that. Define \code{A} once, write each
field as an expression in \code{A}, and change \code{A} to re-evaluate the whole network.
This is how the example projects reproduce both the $A=80$ solution and the $A=0$ answer
printed in the statements.

\section{Three layers}

Values are evaluated in this order, each layer able to use the names of the layers
before it:

\begin{enumerate}
  \item \textbf{Script variables}, the numeric globals of a Python file you link to the
    project (\cref{sec:script}).
  \item \textbf{Project variables}, ordered \emph{name = expression} pairs edited in the
    \gui{Variables} panel. A variable may use the variables above it and the script's.
  \item \textbf{Bindings}, \emph{COMPONENT.field = expression}. They are written into the
    network's fields whenever a variable changes.
\end{enumerate}

The network always holds evaluated numbers. The expressions live beside it, in the
project file's \code{parameters} section, so the studies never see an expression.

\section{The Variables panel}

\begin{minipage}[t]{0.53\linewidth}
The table lists every variable with its \gui{Name}, \gui{Expression}, evaluated
\gui{Value} and \gui{Source}: \code{project} for those defined here, \code{script} for
those read from the linked file.

\begin{itemize}
  \item \gui{Add variable} adds a row. Type its name and expression in the table.
  \item Edit a project variable's name or expression by double-clicking its cell. Script
    variables are read-only here; edit them in the script.
  \item \gui{Remove} deletes the selected project variable.
  \item A variable whose expression fails shows the error as the tooltip of its cell.
\end{itemize}

Every change re-evaluates the bindings, makes the results stale and can be undone. Undo
puts the network back to its previous values too.
\end{minipage}\hfill
\begin{minipage}[t]{0.42\linewidth}\vspace{0pt}
  \centering\shot[0.95\linewidth]{panel-variables.png}
\end{minipage}

\section{Binding a field}

Type an expression instead of a number into any parametric numeric field, in the
property editor or in the data manager. The field turns amber and shows its value
underneath (\cref{sec:properties}). Clear the field to remove the binding. The field then
keeps its last value, or becomes empty if empty is a value it accepts. Bound cells are tinted in the data manager too. Fields marked
\emph{Numbers only} in the reference tables, such as the number of parallel lines, do not
take expressions.

\section{Expressions}

An expression may use:

\begin{itemize}
  \item numbers, including complex ones: \code{1.6+0.3j}. Complex fields also accept
    MATLAB's \code{1.6+0.3i};
  \item the operators \code{+ - * / // \% **} and parentheses;
  \item the variables defined before it;
  \item the functions \code{sqrt}, \code{sin}, \code{cos}, \code{tan}, \code{asin},
    \code{acos}, \code{atan}, \code{atan2}, \code{exp}, \code{log}, \code{log10},
    \code{hypot}, \code{degrees} (also \code{deg}), \code{radians} (also \code{rad}),
    \code{abs}, \code{min}, \code{max}, \code{round}, \code{float} and \code{int};
  \item the constants \code{pi} and \code{e}.
\end{itemize}

Attribute access, indexing, keyword arguments and lambdas are not allowed: expressions
are evaluated by a whitelist, never by Python's \code{eval}. A failing expression says
why, for example \emph{``Unknown name 'B'. Define it as a variable or load it from your
script.''} or \emph{``Division by zero.''}

For example, the external grid of S2 (70~kV, 25~kA at $\cos\varphi = 0.15$) can be
written directly as its Thevenin impedance:

\begin{terminal}
70/(sqrt(3)*25)*(0.15 + 1j*sqrt(1 - 0.15**2))      # = 0.2425 + 1.5983j ohm
\end{terminal}

The example projects use the \gui{Impedance from}~=~\code{ik} mode instead, which says
the same with less algebra.

\section{A variable script}\label{sec:script}

For many variables, or variables computed with some logic, write them in a Python file
of your own, such as \file{examples/variables\_S1.py}, and link it with \gui{Link
script...} in the \gui{Variables} panel. \app\ runs the file and keeps its public numeric
globals: numbers whose names do not start with an underscore.

\begin{caution}
The script is ordinary Python and runs with your privileges. \app\ asks for confirmation
when you link one, and again every time you open a project that links one. If you
decline, the values saved in the project are used. \gui{Reload} runs the script again
when you have edited it. Only link files you trust.
\end{caution}

The script's path is saved in the project. A relative path is read from the project's
folder, which keeps a project and its script movable together. Linking a script is an
undoable step; reloading it is not.

\section{Variables and study cases}

Study cases can override fields with expressions too, and those expressions see the
project variables plus \code{value}, the field's own value (\cref{sec:variant}). A case
that sets \code{LD4.p\_mw} to \code{1.1*value} raises that load by 10\,\% whatever
\code{A} is. On the command line, \cli{--set A=0} changes a variable before solving
(\cref{ch:console}).
